\chapter{Danh mục chi tiết 34 đặc trưng NHANES}
\label{app:feature-inventory}

Bảng \ref{tab:feature-inventory} trình bày đầy đủ danh mục 34 biến số đầu vào trong không gian đặc trưng $F_4$, bao gồm mã định danh của CDC NHANES, tên biến, kiểu dữ liệu và mô tả lâm sàng tương ứng.

\begingroup
\setlength{\tabcolsep}{3pt}
\begin{longtable}{>{\raggedright\arraybackslash}p{2.9cm} >{\raggedright\arraybackslash}p{2.6cm} >{\raggedright\arraybackslash}p{2.5cm} >{\raggedright\arraybackslash}p{6.3cm}}
\caption{Danh mục hệ thống 34 đặc trưng trong không gian $F_4$ theo giao thức chính thức} \label{tab:feature-inventory} \\
\toprule
\textbf{Mã biến} & \textbf{Nhóm đặc trưng} & \textbf{Kiểu dữ liệu} & \textbf{Mô tả lâm sàng / ý nghĩa} \\
\midrule
\endfirsthead
\multicolumn{4}{c}{\small\textbf{Bảng \thetable\ (tiếp theo). Danh mục hệ thống 34 đặc trưng trong không gian $F_4$}} \\
\addlinespace[0.2cm]
\toprule
\textbf{Mã biến} & \textbf{Nhóm đặc trưng} & \textbf{Kiểu dữ liệu} & \textbf{Mô tả lâm sàng / ý nghĩa} \\
\midrule
\endhead
\bottomrule
\endfoot
\multicolumn{4}{l}{\textbf{Nhóm $F_1$ (10 đặc trưng): Dấu hiệu sinh tồn, Nhân trắc \& Xét nghiệm máu}} \\
BMXBMI & Thể chất ($F_1$) & Liên tục (kg/m$^2$) & Chỉ số khối cơ thể (Body Mass Index) \\
BMXWAIST & Thể chất ($F_1$) & Liên tục (cm) & Chu vi vòng eo (đo lường béo phì trung tâm) \\
BPXPLS & Sinh tồn ($F_1$) & Rời rạc (nhịp/phút) & Tần số mạch đập khi nghỉ ngơi (60 giây) \\
SBP\_mean & Sinh tồn ($F_1$) & Liên tục (mmHg) & Huyết áp tâm thu trung bình qua các lần đo \\
DBP\_mean & Sinh tồn ($F_1$) & Liên tục (mmHg) & Huyết áp tâm trương trung bình qua các lần đo \\
LBXGH & Cận lâm sàng ($F_1$) & Liên tục (\%) & Glycohemoglobin (HbA1c - chuyển hóa đường) \\
LBXTC & Cận lâm sàng ($F_1$) & Liên tục (mg/dL) & Cholesterol toàn phần huyết thanh \\
LBDHDD & Cận lâm sàng ($F_1$) & Liên tục (mg/dL) & HDL-Cholesterol (cholesterol tỷ trọng cao) \\
LBXSCR & Cận lâm sàng ($F_1$) & Liên tục (mg/dL) & Creatinine huyết thanh (chức năng lọc của thận) \\
LBXSAL & Cận lâm sàng ($F_1$) & Liên tục (g/dL) & Albumin huyết thanh (chức năng gan, dinh dưỡng) \\
\midrule
\multicolumn{4}{l}{\textbf{Nhóm $F_2$ (19 đặc trưng): $F_1$ + 9 tiền sử bệnh lý lâm sàng}} \\
BPQ020 & Tiền sử ($F_2$) & Phân loại nhị phân & Bác sĩ từng chẩn đoán Tăng huyết áp \\
DIQ010 & Tiền sử ($F_2$) & Phân loại nhị phân & Bác sĩ từng chẩn đoán Đái tháo đường \\
MCQ160B & Tiền sử ($F_2$) & Phân loại nhị phân & Bác sĩ từng chẩn đoán Suy tim sung huyết \\
MCQ160C & Tiền sử ($F_2$) & Phân loại nhị phân & Bác sĩ từng chẩn đoán Bệnh tim mạch vành \\
MCQ160D & Tiền sử ($F_2$) & Phân loại nhị phân & Bác sĩ từng chẩn đoán Đau thắt ngực \\
MCQ160E & Tiền sử ($F_2$) & Phân loại nhị phân & Tiền sử từng bị Cơn đau tim / Nhồi máu cơ tim \\
MCQ160F & Tiền sử ($F_2$) & Phân loại nhị phân & Tiền sử từng bị Đột quỵ / Tai biến mạch máu não \\
MCQ220 & Tiền sử ($F_2$) & Phân loại nhị phân & Tiền sử từng được chẩn đoán mắc Ung thư \\
HUQ010 & Lâm sàng ($F_2$) & Thứ bậc (1--5) & Tự đánh giá tình trạng sức khỏe tổng quát \\
\midrule
\multicolumn{4}{l}{\textbf{Nhóm $F_3$ (29 đặc trưng): $F_2$ + 10 hành vi lối sống \& giấc ngủ}} \\
PAQ605, PAQ620 & Lối sống ($F_3$) & Phân loại nhị phân & Hoạt động thể lực cường độ mạnh / vừa tại nơi làm việc \\
PAQ650, PAQ665 & Lối sống ($F_3$) & Phân loại nhị phân & Hoạt động thể lực cường độ mạnh / vừa trong giải trí \\
PAD680 & Lối sống ($F_3$) & Liên tục (phút/ngày) & Thời gian ngồi tĩnh tại trung bình mỗi ngày \\
smoking\_status & Lối sống ($F_3$) & Phân loại 3 mức & Tình trạng hút thuốc lá (Never, Former, Current) \\
ALQ101 & Lối sống ($F_3$) & Phân loại nhị phân & Uống ít nhất 12 đồ uống có cồn trong năm qua \\
alcohol\_status & Lối sống ($F_3$) & Phân loại 3 mức & Tình trạng dùng rượu bia (Non, Moderate, Heavy) \\
alcohol\_qty & Lối sống ($F_3$) & Thứ bậc lượng hóa & Mức độ tiêu thụ rượu bia lượng hóa \\
SLD010H & Lối sống ($F_3$) & Liên tục (giờ) & Thời lượng giấc ngủ thông thường mỗi đêm \\
\midrule
\multicolumn{4}{l}{\textbf{Nhóm $F_4$ (34 đặc trưng): $F_3$ + 5 nhân khẩu học \& kinh tế xã hội}} \\
RIDAGEYR & Nhân khẩu ($F_4$) & Liên tục (tuổi) & Tuổi của người tham gia khảo sát tại thời điểm khám \\
RIAGENDR & Nhân khẩu ($F_4$) & Phân loại nhị phân & Giới tính sinh học (Nam / Nữ) \\
RIDRETH1 & Nhân khẩu ($F_4$) & Phân loại 5 nhóm & Chủng tộc / Dân tộc tự xác định \\
DMDEDUC2 & Nhân khẩu ($F_4$) & Thứ bậc (1--5) & Trình độ học vấn cao nhất đạt được (người lớn) \\
INDFMPIR & Kinh tế ($F_4$) & Liên tục (tỷ số) & Tỷ số thu nhập gia đình trên chuẩn nghèo liên bang \\
\end{longtable}
\endgroup


\chapter{Cấu hình siêu tham số mô hình}
\label{app:hyperparameters}

\section{Cấu hình Random Forest cố định}
Mô hình cơ sở chính của nghiên cứu sử dụng thuật toán Random Forest với cấu hình siêu tham số xác định từ quy trình Nested Cross-Validation:
\begin{itemize}
    \item Số lượng cây quyết định (\textnormal{n\_estimators}): 150 cây.
    \item Độ sâu tối đa của cây (\textnormal{max\_depth}): 8 tầng.
    \item Tỷ lệ đặc trưng lấy mẫu ngẫu nhiên (\textnormal{max\_features}): 0.5 (50\% số đặc trưng).
    \item Số lượng mẫu tối thiểu để phân nhánh (\textnormal{min\_samples\_split}): 10 mẫu.
    \item Số lượng mẫu tối thiểu tại nút lá (\textnormal{min\_samples\_leaf}): 8 mẫu.
    \item Cơ chế trọng số lớp (\textnormal{class\_weight}): \textnormal{'balanced\_subsample'}.
    \item Tiêu chuẩn chia nhánh (\textnormal{criterion}): Chỉ số Gini (\textnormal{'gini'}).
    \item Hạt giống ngẫu nhiên (\textnormal{random\_state}): Cố định 42.
\end{itemize}

\section{Cấu hình các mô hình cơ sở khác}
\begin{itemize}
    \item \textbf{Logistic Regression (LR):} Chuẩn hóa $L_2$, thuật toán tối ưu lbfgs, số vòng lặp tối đa 1.000, tham số điều hòa $C = 1.0$, \textnormal{class\_weight = 'balanced'}.
    \item \textbf{XGBoost (XGB):} \textnormal{n\_estimators = 100, max\_depth = 4, learning\_rate = 0.05, subsample = 0.8, colsample\_bytree = 0.8}.
    \item \textbf{Ordinal Logistic Regression (OLR):} Tối ưu hóa hàm hợp lý logarit qua SciPy minimize với thuật toán L-BFGS-B, dung sai hội tụ $10^{-6}$, khởi tạo các ngưỡng cắt đơn điệu bằng phân vị mẫu.
\end{itemize}

\FloatBarrier
\chapter{Kết quả chi tiết theo Fold của Nested Cross-Validation}
\label{app:fold-results}
\suppressfloats[t]

Bảng \ref{tab:app-fold-results} tổng hợp chi tiết Macro-F1 trên từng fold kiểm tra ngoài mẫu (Outer Folds) của 4 thuật toán trên tập Legacy Development ($N = 4.045$, bài toán 5 mức $T_5$, không gian đặc trưng $F_4$), trích xuất từ dữ liệu lưu trữ tại \textnormal{nested\_cv\_baseline\_comparison.csv}.

\begin{table}{Chi tiết hiệu năng Macro-F1 theo từng Outer Fold trên không gian $F_4$ ($T_5$)}
\centering
\footnotesize
\setlength{\tabcolsep}{4pt}
\label{tab:app-fold-results}
\begin{tabularx}{\textwidth}{l *{6}{>{\centering\arraybackslash}X}}
\toprule
\textbf{Mô hình} & \textbf{Fold 1} & \textbf{Fold 2} & \textbf{Fold 3} & \textbf{Fold 4} & \textbf{Fold 5} & \textbf{Trung bình $\pm$ SD} \\
\midrule
Logistic Regression & 0.2599 & 0.2198 & 0.2566 & 0.2382 & 0.2292 & 0.2407 $\pm$ 0.0173 \\
Decision Tree       & 0.2538 & 0.2404 & 0.2634 & 0.2239 & 0.2245 & 0.2412 $\pm$ 0.0175 \\
XGBoost (Tuned)     & 0.3017 & 0.2207 & 0.2925 & 0.2341 & 0.2615 & 0.2621 $\pm$ 0.0353 \\
\textbf{Random Forest (Frozen)} & 0.3111 & 0.2615 & 0.2936 & 0.2425 & 0.2690 & \textbf{0.2755 $\pm$ 0.0270} \\
\bottomrule
\end{tabularx}
\end{table}
\FloatBarrier

\chapter{Các biến thể Ma trận Chi phí và Giải thuật CBC}
\label{app:cost-matrices}

\section{Bốn ma trận chi phí Bayes khảo sát trong nghiên cứu}
Nghiên cứu khảo sát 4 ma trận chi phí bất đối xứng $\mathbf{C} \in \mathbb{R}^{3 \times 3}$, trong đó phần tử $\mathbf{C}(i, j)$ định nghĩa chi phí phạt khi mẫu có nhãn thực tế là lớp $i$ nhưng bị mô hình dự đoán nhầm thành lớp $j$ ($i, j \in \{0, 1, 2\}$, tương ứng với $\{C_0, C_1, C_2^+\}$):
\begin{itemize}
    \item \textbf{Ma trận đồng nhất $M_0$ (Uniform Cost / tương đương 0--1 loss):} Chi phí phạt mọi lỗi sai đều bằng 1: $\mathbf{C}(i, j) = 1.0$ với $i \ne j$, và $\mathbf{C}(i, i) = 0$.
    \item \textbf{Ma trận lâm sàng vừa phải $M_1$ (Moderate Penalty):} Phạt lỗi bỏ sót ca nặng sang bình thường $\mathbf{C}(2, 0) = 3.0$; phạt lỗi nhầm nặng sang nhẹ $\mathbf{C}(2, 1) = 1.5$; lỗi nhầm nhẹ sang bình thường $\mathbf{C}(1, 0) = 1.5$; báo động nhầm nặng $\mathbf{C}(0, 2) = 1.5$. Các lỗi phân loại sai khác giữ mức 1.0.
    \item \textbf{Ma trận lâm sàng tăng cường $M_2$ (High Penalty):} Tăng mức phạt bảo vệ ca nặng với $\mathbf{C}(2, 0) = 5.0$; $\mathbf{C}(2, 1) = 2.0$; $\mathbf{C}(1, 0) = 2.0$; $\mathbf{C}(0, 2) = 2.0$.
    \item \textbf{Ma trận lâm sàng tối đa $M_3$ (Extreme Penalty):} Ưu tiên cao nhất cho việc tránh bỏ sót ca nặng với $\mathbf{C}(2, 0) = 8.0$; $\mathbf{C}(2, 1) = 3.0$; $\mathbf{C}(1, 0) = 2.5$; $\mathbf{C}(0, 2) = 2.5$.
\end{itemize}

\section{Giải thuật hiệu chỉnh tính đơn điệu trong Cumulative Binary Classification}
Trong giải thuật phân loại nhị phân tích lũy ($E_{7b}$), hai bộ phân loại nhị phân độc lập dự đoán xác suất vượt ngưỡng:
\begin{align}
    \hat{P}_1 &= P(Y \ge 1 \mid \mathbf{x}) \\
    \hat{P}_2 &= P(Y \ge 2 \mid \mathbf{x})
\end{align}
Về mặt lý thuyết xác suất thứ bậc, điều kiện đơn điệu đòi hỏi $\hat{P}_1 \ge \hat{P}_2$ với mọi $\mathbf{x}$. Tuy nhiên, do hai mô hình được huấn luyện riêng biệt, hiện tượng vi phạm đơn điệu $\hat{P}_2 > \hat{P}_1$ xuất hiện trên thực nghiệm (chiếm 19.53\% ở mô hình Logistic Regression và 5.32\% ở mô hình Random Forest). 

Giải thuật hiệu chỉnh đơn điệu (Isotonic Probability Clipping) được áp dụng khi xuất hiện vi phạm:
\begin{equation}
    \tilde{P}_2 = \min\left(\hat{P}_1, \hat{P}_2\right)
\end{equation}
Khi đó, xác suất suy luận của lớp trung gian $P(C_1 \mid \mathbf{x}) = \hat{P}_1 - \tilde{P}_2 \ge 0$ được bảo đảm không âm về mặt giải tích. Tuy nhiên, tại các điểm xảy ra vi phạm, phép cắt này gán $\tilde{P}_2 = \hat{P}_1$, dẫn đến $P(C_1 \mid \mathbf{x}) = 0$, gây ra hiện tượng suy giảm độ nhạy của lớp trung gian $C_1$ (Recall $C_1 \approx 0\%$) như quan sát được trên thực nghiệm.
